\subsection{Bochner--Riesz 乘子}
\label{sec:ch8-bochner-riesz-multipliers}

如第~\ref{sec:ch1-fourier-integral-convergence-summability}~节所述, Fourier 积分的球形收敛研究: 对 $f\in L^p$, 当 $R\to\infty$ 时
\[
 S_Rf(x)=\int_{|\xi|\le R}\widehat f(\xi)e^{2\pi i\xi\cdot x}\,d\xi
\]
何时在 $L^p$ 范数中或几乎处处逐点收敛到 $f$. 该部分和算子的乘子为 $\chi_{\{|\xi|\le R\}}$:
\[
 \widehat{S_Rf}(\xi)=\chi_{\{|\xi|\le R\}}\widehat f(\xi).
\]
当 $n=1$ 时, 命题~\ref{prop:ch3-interval-multiplier-bound} 已证明该乘子在 $L^p(\mathbb R)$, $1<p<\infty$, 上有界. 当 $n\ge2$ 时, 该乘子性质很差, 见下文第~\ref{subsec:ch8-ball-bochner-riesz}~节. 在高维中, 若考虑比球示性函数更光滑的乘子, 问题较易处理. 例如, 对 $0<t\le R$ 的算子 $S_t$ 取 Cesàro 平均, 得
\[
 \sigma_Rf(x)=\frac1R\int_0^RS_tf(x)\,dt
 =\int_{|\xi|\le R}\left(1-\frac{|\xi|}{R}\right)\widehat f(\xi)e^{2\pi i\xi\cdot x}\,d\xi.
\]
这引出算子族
\[
 \widehat{S_R^af}(\xi)=\left(1-\frac{|\xi|^2}{R^2}\right)_+^a\widehat f(\xi),
 \qquad a>0,
\]
其中 $A_+=\max(A,0)$. 对固定 $a$, $\{S_R^a\}$ 都是 $S_1^a$ 的伸缩, 所以证明其在 $L^p$ 上一致有界只需考虑 $R=1$. 当 $n=1$ 时, $m(\xi)=(1-|\xi|^2)_+^a$ 的 Fourier 变换属于 $L^1$. 事实上, 可分部积分, 再像 Riemann--Lebesgue 引理~\ref{lem:ch1-riemann-lebesgue} 的证明那样得到 $|\widehat m(x)|\le C(1+|x|)^{-1-a}$. 因而 Young 不等式给出 $m$ 在 $L^p$, $1\le p\le\infty$, 上的乘子有界性. 高维则不再成立.

通常用 Bochner--Riesz 乘子取代上述算子:
\[
 \widehat{T^af}(\xi)=(1-|\xi|^2)_+^a\widehat f(\xi),
 \qquad a>0.
\]
取 $\psi_1,\psi_2\in C^\infty$, 支于 $\mathbb R^+$, 且在 $|\xi|\le1$ 时 $\psi_1(|\xi|)=\psi_2(|\xi|)=1$, 则
\[
 (1-|\xi|)_+^a=(1-|\xi|)_+^a[1+|\xi|]^a\psi_1(|\xi|),
\]
\[
 (1-|\xi|^2)_+^a=(1-|\xi|)_+^a[1+|\xi|]^a\psi_2(|\xi|).
\]
方括号中的每个因子都是 Hörmander 乘子, 所以其中一个 $T^a$ 在某个 $L^p$, $1<p<\infty$, 上有界当且仅当另一个有界. 主要结果如下.

\hypertarget{chapter-08-page-181}{}
\begin{Theorem}
\label{thm:ch8-bochner-riesz}
Bochner--Riesz 乘子 $T^a$ 满足:
\begin{enumerate}[label=(\arabic*)]
\item 若 $a>(n-1)/2$, 则 $T^a$ 在 $L^p(\mathbb R^n)$, $1\le p\le\infty$, 上有界.
\item 若 $0<a<(n-1)/2$, 则当
\[
 \left|\frac1p-\frac12\right|<\frac a{n-1}
\]
时 $T^a$ 在 $L^p(\mathbb R^n)$ 上有界; 当
\[
 \left|\frac1p-\frac12\right|\ge\frac{2a+1}{2n}
\]
时 $T^a$ 不有界.
\end{enumerate}
\end{Theorem}

未由定理~\ref{thm:ch8-bochner-riesz} 覆盖的 $p$ 值见下文第~\ref{subsec:ch8-ball-bochner-riesz}~节. $a=(n-1)/2$ 称为临界指标.

Bochner--Riesz 乘子的奇性位于圆周 $|\xi|=1$ 上. 因此把乘子分解到与单位球面距离近似为二进量的二进环带上. 取非负的 $\phi_k\in C^\infty(\mathbb R)$, 使其支于
\[
 1-2^{-k+1}<t<1-2^{-k-1},
\]
对所有多重指标 $\beta$ 满足 $|D^\beta\phi_k|\le C_\beta2^{k|\beta|}$, 且当 $1/2\le t\le1$ 时
\[
 \sum_{k=1}^\infty\phi_k(t)=1.
\]
再定义
\[
 \phi_0(t)=1-\sum_{k=1}^\infty\phi_k(t),
 \qquad0\le t<1/2,
\]
并在 $t\ge1/2$ 时令 $\phi_0(t)=0$. 则
\[
 (1-|\xi|^2)_+^a=\sum_{k=0}^\infty(1-|\xi|^2)^a\phi_k(|\xi|).
\]

\hypertarget{chapter-08-page-182}{}
在 $\phi_k$ 的支集上, $1-|\xi|^2$ 的大小近似 $2^{-k}$, 因而定义
\[
 \widetilde\phi_k(|\xi|)=2^{ka}(1-|\xi|^2)^a\phi_k(|\xi|),
\]
于是
\[
 (1-|\xi|^2)_+^a=\sum_{k=0}^\infty2^{-ka}\widetilde\phi_k(|\xi|).
\]
这把 $T^a$ 分解为
\[
 T^af=\sum_{k=0}^\infty2^{-ka}T_kf,
 \qquad
 \widehat{T_kf}(\xi)=\widetilde\phi_k(|\xi|)\widehat f(\xi).
\]
下一引理描述 $T_k$ 一类算子的行为, 并将给出定理~\ref{thm:ch8-bochner-riesz} 的正面结论.

\begin{Lemma}
\label{lem:ch8-thin-annulus-multiplier}
给定 $0<\delta<1$, 设 $\phi$ 是实直线上的函数, 支于 $1-4\delta<t<1-\delta$, 满足 $0\le\phi\le1$ 和 $|D^\beta\phi|\le C\delta^{-|\beta|}$. 则对任意 $\varepsilon>0$, 与乘子 $\phi(|\xi|)$ 相伴的算子 $T_\delta$ 满足
\begin{equation}
\label{eq:ch8-thin-annulus-estimate}
 \|T_\delta f\|_p
 \le C_\varepsilon
 \delta^{-((n-1)/2+\varepsilon)|2/p-1|}\|f\|_p.
\end{equation}
\end{Lemma}

\begin{Proof}
固定 $K$, 使 $\widehat K(\xi)=\phi(|\xi|)$, 并先设 $b$ 是正偶整数. 由 Fourier 变换的微分恒等式和 Plancherel 定理,
\[
 \|(1+|x|^b)K\|_2
 \le C\|(I+(-\Delta)^{b/2})\phi\|_2
 \le C\delta^{1/2}(1+\delta^{-b})
 \le C\delta^{-b+1/2}.
\]
第二个不等式来自 $\phi$ 支集和 $D^b\phi$ 的大小. 用 Hölder 不等式可把该结论推广到任意 $b>0$. 给定 $b$, 取 $s>1$ 使 $bs$ 为偶整数, 则
\[
 \|(1+|x|^b)K\|_2
 \le C\|(1+|x|^{bs})^{1/s}K\|_2
 \le\|(1+|x|^{bs})K\|_2^{1/s}\|K\|_2^{1/s'},
\]
所需估计由前述论证得到. 取 $b=n/2+\varepsilon$, 再由 Hölder 不等式,
\[
 \|K\|_1
 \le\|(1+|x|^b)K\|_2\|(1+|x|^b)^{-1}\|_2
 \le C_\varepsilon\delta^{-((n-1)/2+\varepsilon)}.
\]

\hypertarget{chapter-08-page-183}{}
因而 Young 不等式在 $q=1,\infty$ 时给出
\[
 \|T_\delta f\|_q
 \le C_\varepsilon\delta^{-((n-1)/2+\varepsilon)}\|f\|_q.
\]
在这两个不等式与 $q=2$ 时的平凡不等式之间插值, 即得 \eqref{eq:ch8-thin-annulus-estimate}.
\end{Proof}

为证明定理~\ref{thm:ch8-bochner-riesz} 的否定部分, 还需要两个引理.

\begin{Lemma}
\label{lem:ch8-compact-multiplier-kernel-lp}
若紧支函数 $m$ 是某个 $L^p$ 上的乘子, 则 $\check m\in L^p$.
\end{Lemma}

\begin{Proof}
取 $f\in\mathcal S$, 使 $\widehat f=1$ 在 $m$ 的支集上成立. 则 $f\in L^p$ 且 $T_mf\in L^p$, 而 $\widehat{T_mf}=m\widehat f=m$.
\end{Proof}

\begin{Lemma}
\label{lem:ch8-bochner-riesz-kernel}
$(1-|\xi|^2)_+^a$ 的 Fourier 变换为
\begin{equation}
\label{eq:ch8-bochner-riesz-kernel}
 K^a(x)=\pi^{-a}\Gamma(a+1)|x|^{-n/2-a}J_{n/2+a}(2\pi|x|),
\end{equation}
其中 $J_\mu$ 是 Bessel 函数
\[
 J_\mu(t)=\frac{(t/2)^\mu}{\Gamma(\mu+1/2)\Gamma(1/2)}
 \int_{-1}^1e^{its}(1-s^2)^{\mu-1/2}\,ds.
\]
\end{Lemma}

\begin{Proof}
公式 \eqref{eq:ch8-bochner-riesz-kernel} 立即来自以下两个结果:
\begin{enumerate}[label=(\arabic*)]
\item 若径向函数 $f(x)=f_0(|x|)\in L^1(\mathbb R^n)$, 则 $\widehat f$ 径向, 且
\[
 \widehat f(\xi)=2\pi|\xi|^{1-n/2}
 \int_0^\infty f_0(s)J_{n/2-1}(2\pi|\xi|s)s^{n/2}\,ds.
\]
\item 若 $j>-1/2$, $k>-1$, $t>0$, 则
\[
 J_{j+k+1}(t)=\frac{t^{k+1}}{2^k\Gamma(k+1)}
 \int_0^1J_j(ts)s^{j+1}(1-s^2)^k\,ds.
\]
\end{enumerate}
证明见 Stein 和 Weiss \MBUnresolvedReference{citation}{[18, pp. 155, 170]}.
\end{Proof}

\begin{Proof}[定理~\ref{thm:ch8-bochner-riesz} 的证明]
由引理~\ref{lem:ch8-thin-annulus-multiplier}, 对每个 $k$,
\[
 \|T_kf\|_p
 \le C_\varepsilon2^{k((n-1)/2+\varepsilon)|2/p-1|}\|f\|_p.
\]
所以 Minkowski 不等式给出
\[
 \|T^af\|_p
 \le C_\varepsilon\sum_{k=0}^\infty
 2^{k((n-1)/2+\varepsilon)|2/p-1|-ka}\|f\|_p.
\]
正面结论立即得到.

\hypertarget{chapter-08-page-184}{}
为得否定结论, 从 Bessel 函数的定义可知 $J_a(t)=O(t^a)$, $t\to0$. 又当 $t\to\infty$ 时, $J_a(t)\approx t^{-1/2}$, 见 Stein \MBUnresolvedReference{citation}{[17, p. 338]}. 因而引理~\ref{lem:ch8-bochner-riesz-kernel} 给出
\[
 |K^a(x)|\le C\quad\text{当 }|x|\to0,
\]
以及
\[
 |K^a(x)|\approx|x|^{-(n+1)/2-a}quad\text{当 }|x|\to\infty.
\]
所以 $K^a\in L^p$ 当且仅当
\[
 p>\frac{2n}{n+1+2a}.
\]
结合引理~\ref{lem:ch8-compact-multiplier-kernel-lp} 与对偶论证, 得当
\[
 \left|\frac1p-\frac12\right|\ge\frac{2a+1}{2n}
\]
时 $T^a$ 不有界.
\end{Proof}
